\section{A logarithmic Hodge vector and its coefficient lattice}
\label{sec:logarithmic-input}

Let $f:(X,E)\to(Y,D)$ satisfy the hypotheses of
\Cref{thm:upper-bound}, and put $d=\dim X-\dim Y$ and $L_Y=K_Y+D$.
Fix one total-space logarithmic pluriform
\begin{equation}\label{eq:log-input-section}
  0\ne s\in H^0\bigl(X,\OO_X(m(K_X+E))\bigr),\qquad m>0.
\end{equation}
We turn $s$ into a highest Hodge vector whose coefficients retain
the original base cotangent lattice. The open set on which a cyclic
cover constructed from $s$ is a smooth family will usually be smaller
than $V=Y\setminus D$. The differential lattice nevertheless remains
$\Omega_Y^1(\log D)$, including at divisors in $V$ where that cover
degenerates. Retaining this lattice will later turn the Hodge vector
into actual base sections.

\phantomsection\label{sec:preliminaries}
\subsection{Variations and extension conventions}

A polarized rational variation of pure Hodge structure of weight $w$ on a
smooth quasi-projective variety $B^\circ$ consists of a rational local
system $\mathbb V_{\Q}$, its flat holomorphic bundle
$\mathcal V=\mathbb V_{\Q}\otimes_{\Q}\OO_{B^\circ}$ with connection
$\nabla$, a decreasing holomorphic filtration $F^\bullet\mathcal V$,
and a flat polarization.  The filtration defines polarized pure Hodge
structures on fibers and satisfies Griffiths transversality
\[
 \nabla(F^p\mathcal V)\subseteq
 F^{p-1}\mathcal V\otimes\Omega^1_{B^\circ}.
\]
The initial variations in this paper have invariant lattices; torsion is
discarded when taking such a lattice.  A complex variation used below is
obtained, after a specified finite level if necessary, as a polarized
complex Hodge summand of a tensor construction in these rational
variations and their duals.  In particular its projector is flat and
preserves the Hodge decomposition.  Such a summand need not have its own
real form.  It has the induced positive Hodge metric and a flat Hermitian
polarization.  Constant changes of the Hodge indices in these constructions
do not change the highest Hodge spaces or their differentials.

We record precisely the standard extension results used here.  Choose a
smooth projective compactification $B$ with SNC boundary
$\Delta=B\setminus B^\circ$.  The local boundary monodromies of the
integral polarized variation are quasi-unipotent, by the monodromy theorem
\cite[(6.1), pp.~245--246]{Schmid}.  They become unipotent on a finite normal level.  One
can choose the kernel of reduction of the lattice modulo an integer
$M\geq3$: if a quasi-unipotent integral matrix $T$ is congruent to $1$
modulo $M$, then for every eigenvalue $\zeta$ the number
$(\zeta-1)/M$ is an algebraic integer.  All its conjugates have absolute
value at most $2/M<1$, so its algebraic norm forces it to be zero.
Thus every eigenvalue of $T$ is $1$.  The cover exists algebraically by
Riemann existence; normalize its compactification and resolve it.  New
boundary monodromies are products of powers of commuting unipotent
monodromies and remain unipotent.  This level construction may be combined
with a finite level killing a monodromy component group.

At unipotent level, let $\overline{\mathcal V}$ denote Deligne's canonical
extension, with nilpotent logarithmic residues.  The nilpotent orbit
theorem extends $F^\bullet$ by holomorphic subbundles of
$\overline{\mathcal V}$; the induced logarithmic Higgs field is
\[
 \theta:\Gr_F^p\overline{\mathcal V}\longrightarrow
 \Gr_F^{p-1}\overline{\mathcal V}\otimes\Omega^1_B(\log\Delta).
\]
We use the nilpotent orbit and limiting mixed Hodge structure theorems
of \cite[(4.9), (4.12), and (6.16)]{Schmid}.

Here and later, pullback compatibility means compatibility at unipotent
level.  Suppose $g:(B',\Delta')\to(B,\Delta)$ is a morphism of smooth
compactifications whose open part maps to $B^\circ$, and that
$g^{-1}\Delta$ is supported on the SNC divisor $\Delta'$.  Locally write
\[
 g^*z_i=v_i\prod_j t_j^{a_{ij}},\qquad v_i\ \text{a holomorphic unit}.
\]
If $N_i$ are the commuting nilpotent residue matrices, the residues of the
pulled-back logarithmic connection are $\sum_i a_{ij}N_i$ and are
nilpotent.  Uniqueness of the canonical extension therefore identifies it
with $g^*\overline{\mathcal V}$.  The pulled-back Hodge subbundles are
subbundles and agree with the Hodge filtration on the open set, so
uniqueness of the extending Grassmann maps identifies the extended
filtrations as well.  The units $v_i$ contribute only invertible
holomorphic exponential factors.  This proves compatibility also under
the modifications and ramified covers used below.  Tensor operations and
flat Hodge projectors commute with these extensions: the latter commute
with the local monodromies and hence with the logarithmic exponentials.
Thus all the preceding extension assertions apply to the
complex Hodge summands specified above.

The regular-singular Riemann--Hilbert correspondence
\cite[Chapter II, \S\S5--6]{DeligneRSE} makes the flat bundles, their
tensor morphisms, and their pullbacks algebraic on the open set.  On a
projective compactification the extending Hodge subbundles are algebraic
by GAGA; this gives algebraicity on the original open set as well, using
finite descent when a level was introduced.  We also use the theorem of
the fixed part: the monodromy-invariant subspace of a polarized variation
on a smooth quasi-projective base is a constant Hodge substructure, and
its inclusion is a morphism of variations.  This applies to tensor
constructions and finite étale covers
\cite[\S4.1]{DeligneHodgeII} and \cite[(7.22)]{Schmid}.
The assertions about connected algebraic
monodromy needed for the marked-period construction will be stated at
their use.

\subsection{The vector and the coefficient statement}
We now specify regularity in the logarithmic cotangent lattice on a
higher model.
Let $\beta:Y'\to Y$ be a projective birational morphism with $Y'$ smooth,
and let $D'$ be the reduced union of the strict transform of $D$ and all
$\beta$-exceptional divisors, with simple normal crossings.  Set
\[
 L_{Y'}=K_{Y'}+D',\qquad B'=\Omega^1_{Y'}(\log D').
\]
All rational identifications of canonical lines use the canonical
pullback of rational differential forms.  At the generic point of a
prime $A\subset Y'$, choose local frames of $\OO_{Y'}(-aL_{Y'})$ and
$(B')^{\otimes j}$.  After a sufficiently divisible transverse
ramification, a variation with quasi-unipotent monodromy has unipotent
monodromy.  A tensor has \emph{regular coefficients} if, in the
Schmid extension of its Hodge factor, its coefficients relative to the
\emph{pulled-back chosen frames} are regular.  Thus this convention
pulls back the vector bundle $(B')^{\otimes j}$; it does not replace
that bundle by cotangent tensors on the testing curve.

\begin{proposition}[Logarithmic Hodge input]\label{prop:logarithmic-vector}
There are a dense smooth open $Y^\circ\subset V$, a polarizable pure
rational variation $\mathbb V$ on $Y^\circ$ with an invariant lattice
and quasi-unipotent boundary monodromy, and a nonzero algebraic map
\begin{equation}\label{eq:log-input-vector}
 u:\OO_{Y^\circ}(-L_Y)\longrightarrow
 \mathcal E^d,
 \qquad
 \mathcal E^p=\Gr_F^p(\mathbb V\otimes\OO_{Y^\circ}),
 \qquad F^{d+1}\mathbb V=0.
\end{equation}
On every model $(Y',D')$ just described, the rationally identified map
and all its Higgs iterates
\begin{equation}\label{eq:log-input-iterates}
 \theta^j(u):\OO_{Y'}(-L_{Y'})\dashrightarrow
 \mathcal E^{d-j}\otimes(B')^{\otimes j}
\end{equation}
have regular coefficients at every prime of $Y'$ in the preceding
sense.  Here the variation and its Hodge bundles are understood over
the common dense open, and a grade outside their range is zero.

The same statement holds for the homogeneous tensor replacement used
in \Cref{prop:marked-period}.  More precisely, if $u_*$ is obtained
from $u$ by homogeneous tensor operations of positive degree $a$ and
flat Hodge-compatible projections, then
\begin{equation}\label{eq:log-input-tensor}
 u_*:\OO(-aL_{Y'})\dashrightarrow E_0,
 \qquad
 \theta^j(u_*):\OO(-aL_{Y'})\dashrightarrow
 E_j\otimes(B')^{\otimes j}
\end{equation}
have the same coefficient property. Here $\mathbb A$ is the resulting
pure variation, $p_*$ is its highest nonzero Hodge index, and
$E_j=\Gr_F^{p_*-j}\mathbb A$. Finite levels and a common
further unipotent ramification may be used in these assertions.
\end{proposition}

\begin{proof}
We give the construction and then prove the assertion for the full
coefficient lattice.

\smallskip\noindent\emph{Step 1: the cyclic root and the pure vector.}
The cyclic-cover construction and its Higgs iteration follow the
Viehweg--Zuo method \cite[\S3]{ViehwegZuoIsotriviality}; here we must also
retain every coefficient in the original base lattice.
Put $N=K_X+E$. Normalize the cyclic cover defined by an $m$-th root
of $s$, and resolve it to obtain a projective generically finite map
$\psi:Z\to X$.  The source may be disconnected.  The tautological
root gives a regular pairing
\begin{equation}\label{eq:log-root-pairing}
 \psi^*\OO_X(-N)\longrightarrow\OO_Z.
\end{equation}
Choose a reduced simple-normal-crossing divisor on $Z$ containing
the inverse image of $E$ and the exceptional divisors needed for the
chosen compactification.  Shrink the base to $Y^\circ\subset V$ so
that $Z\to Y$ and every stratum of this divisor are smooth there.
Write $H$ for its restriction over $Y^\circ$ and
$q:Z\setminus H\to Y^\circ$ for the resulting open family.  Properness
of the compactification and smoothness of all its strata give local
topological triviality of this family.

The logarithmic mixed Hodge construction gives a variation on
$R^d q_*\Q$ with weight filtration $W$ and Hodge filtration $F$.
We recall the precise ingredients used here.  For a smooth projective
compactification with a simple-normal-crossing boundary, residues
identify the weight quotients of the logarithmic complex with
complexes on the closed boundary intersections, shifted and Tate
twisted.  The associated weight spectral sequence degenerates at
$E_2$, the Hodge-to-de Rham sequence degenerates at $E_1$, and the
same residue filtration on a fixed form sheaf degenerates at $E_2$;
see \cite[\S3.1.5, Theorem~3.2.5 and Corollary~3.2.13]{DeligneHodgeII}.
These statements apply relatively on $Y^\circ$: cohomology and base
change give the Hodge bundles, and the proper smooth boundary strata
give the pure Gauss--Manin systems occurring on the residue page.
Their Gysin maps, including the Tate twists, are morphisms of pure
variations.  Their kernels and images split by polarization.  Thus
the weight quotients are polarizable rational variations, and their
integral cohomological constructions give invariant lattices after
discarding torsion.  Their monodromy is quasi-unipotent.  The proper
stratum Gauss--Manin connections are algebraic and regular singular
by the regularity theorem for Gauss--Manin connections
\cite[Chapter~II, \S6.14 and Theorem~7.9]{DeligneRSE}; residues and their cohomological maps
are algebraic.  Consequently all the Hodge and weight maps used
below have their usual algebraic structures.

Over $Y^\circ$ there is a canonical identification
\[
 \Omega^d_{X/Y}(\log E)\otimes\OO_X(-N)
       =f^*\OO_Y(-L_Y).
\]
Pulling back the top relative logarithmic form and applying
\eqref{eq:log-root-pairing} therefore gives
\begin{equation}\label{eq:log-input-mixed-vector}
 u^{\mathrm{mix}}:\OO_{Y^\circ}(-L_Y)
   \longrightarrow F^d(R^d q_*\C\otimes\OO_{Y^\circ}).
\end{equation}
The form is nonzero: the root is nonzero generically, and pullback
of a top relative differential under a generically finite map in
characteristic zero is injective at the generic point.  A nonzero
global top logarithmic form represents a nonzero top Hodge class
by Hodge-to-de Rham degeneration.  Choose the smallest $w$ for which
the image of \eqref{eq:log-input-mixed-vector} lies generically in
$W_w$.  Projection to $\Gr^W_w$ is nonzero.  Let $\mathbb V=\Gr^W_w$
and let $u$ be this projection.  Since the fibers of $q$ have
dimension $d$, $F^{d+1}=0$.  Thus $u$ is in the highest nonzero
Hodge grade of this pure variation, even when it does not span that
grade.

For use with this construction, the Higgs field has the following
description.  Filter absolute logarithmic forms by the number of
base differentials.  The short exact sequence consisting of base
degrees zero and one induces on relative de Rham cohomology the
Gauss--Manin connection.  This follows either from the logarithmic
de Rham comparison or by lifting a flat class to an absolute class
on a topologically trivial base disk.  Taking its Hodge symbol
gives the connecting map for the corresponding short exact sequence
of individual form sheaves.  Hodge-to-de Rham degeneration identifies
the resulting cohomology sheaves with the Hodge grades.  Since $W$
is flat and all the filtrations are strict, these connecting maps
preserve weight and induce the pure Higgs field on $\Gr^W_w$.

\smallskip\noindent\emph{Step 2: coefficient maps on higher models.}
We now work on a higher model $Y'$.  Resolve the main component of
$X\times_Y Y'$ to obtain $f':X'\to Y'$, and let $E'$ be the strict
transform of $E$ together with the reduced divisors exceptional
over $X$.  Resolve the pair, preserving the smooth relative pair
over $Y'\setminus D'$.  Then
\begin{equation}\label{eq:log-input-support}
 \Supp f'^*D'\subset\Supp E'.
\end{equation}
Indeed, over $D$ this follows from the original support inclusion.
Outside $D$, the image in $Y$ of a $\beta$-exceptional divisor has
codimension at least two.  Smoothness of $f$ there implies that its
inverse image has codimension at least two in $X$.  Every divisor
above it is consequently exceptional over $X$ and belongs to $E'$.
The logarithmic discrepancy formula gives a pulled-back section
\[
 s'\in H^0(X',\OO_{X'}(mN')),
 \qquad N'=K_{X'}+E',
\]
whose rational identification with $s$ is fixed by pullback.
We may dominate both cyclic-root diagrams on a common resolution,
preserving the chosen family over a common dense open.

Fix a prime $A\subset Y'$.  On the inverse image of a sufficiently
small Zariski neighborhood of its generic point, the natural map
\begin{equation}\label{eq:log-input-subbundle}
 f'^*B'\longrightarrow\Omega^1_{X'}(\log E')
\end{equation}
is an injection of subbundles.  If $A\not\subset D'$, this is the
relative smooth-pair condition.  For a component of $D'$, take its
parameter $t$ and tangential base coordinates $y_2,\ldots,y_b$.
At a point over a general point of $A$, write
$f'^*t=v\prod z_\alpha^{n_\alpha}$ in boundary coordinates, with
$v$ a unit and at least one $n_\alpha>0$.  Its logarithmic
differential has the nonzero residue vector $(n_\alpha)$.
On each boundary stratum dominating $A$, generic smoothness gives
independent pullbacks of $dy_2,\ldots,dy_b$ in the ordinary
tangential directions.  Strata not dominating $A$, and the
nonsmooth loci on the dominating strata, have images missing a
dense open of $A$; properness permits their removal on the base.
The residue vector and these tangential differentials are linearly
independent.  This proves \eqref{eq:log-input-subbundle} at every
point over that open of $A$, without any smoothness assumption
over the original boundary.

Let $\Omega^i_{\mathrm{rel}}$ denote the exterior powers of the
locally free quotient in \eqref{eq:log-input-subbundle}.  Its top
determinant gives
\begin{equation}\label{eq:log-input-top-identity}
 \Omega^d_{\mathrm{rel}}\otimes\OO_{X'}(-N')
       =f'^*\OO_{Y'}(-L_{Y'}).
\end{equation}
For $i\ge1$, the exterior-power filtration gives the exact sequence
\begin{equation}\label{eq:log-input-absolute-extension}
 \begin{split}
 0\longrightarrow f'^*B'\otimes\Omega^{i-1}_{\mathrm{rel}}
  \longrightarrow
  \Omega^i_{X'}(\log E')/\mathrm{Fil}_{\mathrm{base}}^2
  \longrightarrow\Omega^i_{\mathrm{rel}}
  \longrightarrow0.
 \end{split}
\end{equation}
The following connecting-map construction is the logarithmic cotangent
procedure of Kov\'acs \cite[\S1, Lemma~1.1]{KovacsFineModuli}, applied
before passing to the root family so that the coefficient lattice remains
fixed. Tensor this sequence by $\OO_{X'}(-N')$, apply higher direct
images, and use the projection formula.  Starting with the unit
map into the direct image of \eqref{eq:log-input-top-identity},
its successive connecting maps give honest morphisms of coherent
sheaves
\begin{equation}\label{eq:log-input-coherent-iterates}
 \alpha_j:\OO_{Y'}(-L_{Y'})\longrightarrow
 A_j\otimes(B')^{\otimes j},\qquad
 A_j=R^j f'_*(\Omega^{d-j}_{\mathrm{rel}}\otimes\OO_{X'}(-N')).
\end{equation}
They are defined before passing to the root family.  On the good
open, differential pullback and root multiplication map each
short exact sequence \eqref{eq:log-input-absolute-extension} to
the corresponding sequence on that family.  Naturality of
connecting homomorphisms identifies the images of $\alpha_j$
with $(\theta^{\mathrm{mix}})^j(u^{\mathrm{mix}})$.  This argument
uses morphisms of short exact sequences of $\OO$-modules; it
does not require root multiplication to commute with the de Rham
differential or a connection on $\OO(-N')$.

\smallskip\noindent\emph{Step 3: the semistable coefficient test.}
Take a general algebraic curve transverse to $A$ and then a disk
about its intersection point.  Resolve the root diagram over
that disk, keeping its smooth pair over the puncture, and include
the special fiber in the resolved boundary.  One-parameter
semistable reduction, after finite ramification, gives a
projective semistable family with smooth total space
\[
 g:Z_\Delta\longrightarrow\Delta
\]
with reduced simple-normal-crossing special fiber $V_0$, and a
horizontal divisor $H$ such that $V_0+H$ has simple normal
crossings.  We use the toroidal construction with the horizontal
divisor included in the boundary
\cite[Chapter~IV, \S3]{KKMS}; an explicit statement allowing
the prescribed boundary is \cite[Theorem~7.17]{KollarMori}.
After resolving the special fiber, its parameter is monomial;
the resolution centers over the bad fiber, followed by base
change and toroidal subdivisions for the vertical
monomial preserve the transverse horizontal coordinate
hyperplanes.  Thus the local equations on the resulting model
are
\begin{equation}\label{eq:log-input-semistable-coordinates}
 \tau=z_1\cdots z_r,
 \qquad H:\ z_{r+1}\cdots z_{r+h}=0.
\end{equation}
In particular every closed horizontal intersection is itself a
proper semistable family over $\Delta$.  The model maps to $X'$
and to the chosen root family.  Its boundary contains the
inverse image of $E'$.

Let $i:\Delta\to Y'$ denote the ramified testing map.  Pullback
of forms, followed by the root pairing, gives
\[
 p^*(\Omega^{d-j}_{\mathrm{rel}}\otimes\OO_{X'}(-N'))
 \longrightarrow
 \Omega^{d-j}_{Z_\Delta/\Delta}(\log(V_0+H)),
\]
where $p:Z_\Delta\to X'$ is the diagram map.  This descends to
relative forms because every pulled-back base logarithmic form
from $Y'$ is a form from $\Delta$ and vanishes in the relative
quotient.  Its pole at $0$ is at most logarithmic.  The root
pairing is regular on this model; alternatively, its monic
power equation and normality prove this regularity.
The natural coherent-cohomology pullback maps therefore give
\begin{equation}\label{eq:log-input-coefficient-test}
 i^*A_j\longrightarrow
 R^j g_*\Omega^{d-j}_{Z_\Delta/\Delta}(\log(V_0+H)).
\end{equation}
No cohomology base-change isomorphism at $0$ is needed.

To see exactly what \eqref{eq:log-input-coefficient-test} tests,
choose a local source generator $e$ and an entire local frame
$(b_\nu)_\nu$ of $(B')^{\otimes j}$, and write
\[
 \alpha_j(e)=\sum_\nu a_\nu\otimes b_\nu.
\]
Every $i^*a_\nu$ maps to a regular relative log-form cohomology
class on $\Delta$.  We keep each $i^*b_\nu$ as a basis vector
of the external bundle $i^*(B')^{\otimes j}$.  On $\Delta^*$,
smooth-family comparison identifies these classes with the
individual coefficients of the original mixed Higgs iterate.
Each coefficient belongs to $W_w$ there, because the mixed
Higgs field preserves $W_w$.

\Cref{lem:weight-extension} shows that its weight projection is
a regular section of the Schmid-extended pure Hodge grade.
This proves the required bound for every coefficient of
\eqref{eq:log-input-iterates}.  In particular, a tangential
factor annihilated by the differential of $i$ is still retained
as an external frame vector.  At an additional root-cover
discriminant outside $D'$, the transverse frame is $dt$.
Although $d(\tau^e)=e\tau^e d\log\tau$, no coefficient is
divided by $e\tau^e$, since this differential substitution is
never made on the external frame.  This proves the bound with
the original $D'$, without adding that discriminant to it.

For completeness, these tests detect divisorial regularity.
Take a finite level on which the variation is unipotent and a
smooth compactification with normal-crossing boundary.  Its
canonical Hodge extensions are algebraic, and the rational
coefficient maps are meromorphic in their frames.  A pole at
a prime persists on a general transverse disk; the preceding
argument excludes it after divisible ramification.  Further
ramified tests preserve regularity by compatibility of the
unipotent canonical extension with pullback.  Since $Y'$ and
$A$ were arbitrary, this proves the assertion on every stated
higher model.

\smallskip\noindent\emph{Step 4: homogeneous replacements.}
Finally, the tensor product Higgs field satisfies the Leibniz
rule.  Each coefficient of $\theta^j(u^{\otimes a})$ is a finite
sum of tensor products of coefficients of iterates of $u$,
with total iteration degree $j$.  Tensor products and flat
Hodge-compatible projections commute with unipotent canonical
extension.  The same applies to the conjugate factors in a
finite-orbit tensor construction on a common finite level.
Thus every homogeneous tensor construction used in
\Cref{prop:marked-period} preserves the asserted regularity,
with source $\OO(-aL_{Y'})$.  Those constructions use polynomial
tensor operations and projections, so no division introduces
a pole.  This proves \eqref{eq:log-input-tensor}.
\end{proof}

The preceding construction reduces regularity to one precise extension
fact. It is important that the weight steps extend as subbundles: this
allows a scalar zero to be divided out before taking its pure projection.

\begin{lemma}[Extension of a weight projection]\label{lem:weight-extension}
Suppose $g:Z\to\Delta$ is projective and has the local form
\eqref{eq:log-input-semistable-coordinates}, with $Z$ smooth,
reduced special fiber $V_0$, and horizontal boundary $H$.
For fixed $i,j$, the relative log-form cohomology
\[
 \mathcal A^{i,j}=R^j g_*\Omega^i_{Z/\Delta}(\log(V_0+H))
\]
is locally free and carries a filtration by subbundles extending
the weight filtration on its Hodge grade over $\Delta^*$.
Its successive quotients are the Hodge-graded unipotent
canonical extensions of the corresponding pure weight
variations.  In particular, a regular class that belongs
generically to a weight step projects regularly to that
step's pure weight quotient.
\end{lemma}

\begin{proof}
Filter $\Omega^i_{Z/\Delta}(\log(V_0+H))$ by the number $k$ of
horizontal logarithmic factors.  In
\eqref{eq:log-input-semistable-coordinates}, taking a horizontal
residue commutes with quotienting by
$d\log\tau=\sum_{a=1}^r d\log z_a$.  Therefore the residue
quotient of filtration degree $k$ is
\begin{equation}\label{eq:log-input-relative-residues}
 \bigoplus_{|I|=k}
 (\iota_I)_*\Omega^{i-k}_{H_I/\Delta}(\log(V_0|_{H_I})),
 \qquad H_I=\bigcap_{\alpha\in I}H_\alpha.
\end{equation}
The terms are zero when their form degree is negative.
The maps $g_I:H_I\to\Delta$ are proper projective semistable
families.  The semistable logarithmic de Rham comparison
\cite[\S2, Theorem~2.7, Corollaries~2.9--2.10 and
Theorem~2.11]{Steenbrink77}, in the canonical-extension form explained
in \cite[\S4.9 and the proof of Lemma~4.10, Step~1]{FF14}, with its
strictness argument completed in \cite{FF14Remark}, identifies
\[
 R^q(g_I)_*\Omega^p_{H_I/\Delta}(\log(V_0|_{H_I}))
 \quad\text{with}\quad
 \Gr_F^p\overline{R^{p+q}(g_I|_{\Delta^*})_*\C},
\]
where the bar means unipotent canonical extension with its
extended Hodge filtration.  In particular these are locally
free sheaves.

Consider the spectral sequence of the finite residue filtration
of this single relative form sheaf.  We use the indexing
\[
 E_1^{-k,j+k}=
 \bigoplus_{|I|=k}R^j(g_I)_*
 \Omega^{i-k}_{H_I/\Delta}(\log(V_0|_{H_I})).
\]
On $\Delta^*$, $d_1$ is the Hodge grade of the signed Gysin
map from the $k$-fold intersections to the $(k-1)$-fold
intersections, with the Tate twists included.  Source and
target have the same pure weight $i+j+k$.  A morphism of
polarizable pure variations has a flat Hodge kernel and image
and admits flat Hodge complements.  Consequently it factors
as a projection to a complement, an isomorphism onto its image,
and an inclusion.  Canonical extension and extension of $F$
preserve this direct-sum factorization.  Its extended
Hodge-graded map has constant rank at $0$.

The coherent $d_1$ and that extended Gysin map are holomorphic
maps between the same locally free $E_1$ terms and agree on
$\Delta^*$.  They therefore agree everywhere.  It follows that
$E_2$ is locally free and is precisely the canonical
Hodge-graded extension of the pure residue cohomology.
By the fixed-form version of Deligne's residue theorem
\cite[Corollary~3.2.13(iii)]{DeligneHodgeII}, every $d_r$ for
$r\ge2$ vanishes on $\Delta^*$.  Inductively its target is
locally free on $\Delta$, so a holomorphic map vanishing on
the puncture vanishes everywhere.  Thus $E_2=E_\infty$ on
the disk.  This is also the constant-rank extension mechanism
in \cite[Proposition~3.12 and the proof of Lemma~4.10,
Step~1]{FF14}; here the residue calculation
\eqref{eq:log-input-relative-residues} supplies its particular
geometric complex.

The abutment has a finite filtration with these locally free
successive quotients.  It and every filtration step are
therefore locally free, and the filtration steps are
subbundles.  On $\Delta^*$ this filtration, with the usual
cohomological shift, is the induced weight filtration on
$\Gr_F^i H^{i+j}(Z_t\setminus H_t)$: this is exactly the
compatibility of the two filtrations in the logarithmic mixed
Hodge complex, or equivalently the fixed-form assertion in
the cited corollary.  Its extended quotients were identified
on $E_2$ above.  Finally, the image of a regular class in the
quotient by a weight subbundle is a section of a locally free
sheaf.  If it is generically zero it is zero.  The class hence
lies in that subbundle on the entire disk and has a regular
image in its successive quotient, as claimed.
\end{proof}

We now retain a zero on a specified original fiber. The source line in
this calculation is the pulled-back original line, whereas the preceding
higher-model estimates used the new logarithmic canonical line. Keeping
these frames distinct is what detects vanishing at the chosen point.

\begin{lemma}[Zero test in the original source frame]\label{lem:zero-test}
Suppose $\dim Y>0$ and $y\in V$, and suppose the section $s$ in
\eqref{eq:log-input-section} vanishes identically on $f^{-1}(y)$.
Blow up $y$ on the base if $\dim Y>1$, and let $A$ be its
exceptional divisor; if $\dim Y=1$, put $A=\{y\}$.
Along a general transverse test at $A$, after ramification
of order $e$ divisible by $m$, the coefficients of $u$ in
the pulled-back \emph{original} $\OO_Y(-L_Y)$ frame vanish
to order at least $e/m$.  Every homogeneous replacement
$u_*$ of degree $a>0$ vanishes to order at least $ae/m$
in the corresponding original $\OO_Y(-aL_Y)$ frame.
\end{lemma}

\begin{proof}
Near $y$ the original morphism is smooth and its boundary is
relative simple normal crossing.  Its fiber is reduced.
Consequently, in any local frame of $\OO_X(m(K_X+E))$, the
coefficient of $s$ belongs to
$\mathfrak m_y\OO_X$.  On the blowup the ideal
$\mathfrak m_y$ becomes the invertible ideal of $A$.
Along a transverse parameter $t$ at its general point the
pulled coefficient is therefore divisible by $t$.

Use throughout the pulled-back old line $\OO_X(K_X+E)$ and
the old relative logarithmic forms.  After $t=\tau^e$, the
tautological root $r$ satisfies, in this old line frame,
\begin{equation}\label{eq:log-input-zero-equation}
 \left(r/\tau^{e/m}\right)^m=s/\tau^e.
\end{equation}
The right-hand side is regular.  The left-hand root is a
rational section integral over the normal root model and
is therefore regular there.  It stays regular on the
semistable model used above.  Thus the map from old relative
top logarithmic forms into the semistable relative forms
has the factor $\tau^{e/m}$.

Passing to coherent cohomology retains that factor.  Divide
it out first: the resulting regular class still lies
generically in $W_w$, since multiplication by a nonzero
scalar on the puncture does not change membership in a
weight step.  By \Cref{lem:weight-extension}, its projection
to $\Gr^W_w$ is regular.  Multiplying the projected class
back by $\tau^{e/m}$ proves the claimed order for $u$.
Every degree-$a$ homogeneous tensor operation multiplies
the order by $a$, and every flat Hodge-compatible projection
preserves this divisibility.  The result follows for $u_*$.
This calculation never replaces the old pulled-back source
line by $\OO(-L_{Y'})$ on the blown-up base; the asserted
positive order is in the original source frame.
\end{proof}
